\section{Mathematical formulation}
\label{sec:governing}

The formulation follows the order in which one enthalpy generates the constitutive relations, the field equations, the boundary conditions at the faces and the edge, and the power flux that identifies a backward mode.  The section ends with the dimensionless groups, the reductions the model must reproduce, and the condition for positive definiteness of the elastic--gradient--carrier block.

\subsection{Geometry, kinematics and the gradient energy}
\label{sec:governing_variational}
\label{sec:governing_kinematics}

Consider a transversely isotropic ($6mm$) n-type ZnO plate of thickness $d$ whose faces are the planes $z=\pm d/2$.  The plate occupies $x<0$ and ends at a traction-free edge at $x=0$.  It is unbounded in $y$, so a single tangential wavenumber $k_y$ is conserved and every field has the form $f(x,z)\,e^{\ii k_yy-\ii\omega t}$ with all three displacement components retained.  The formulation is therefore two-and-a-half-dimensional: the $(x,z)$ cross-section is discretized, $y$ is not.  The independent fields, with $\vect x=(x,y,z)$ and the $y$ dependence carried by the factor $\ee^{\ii k_yy}$ so that the amplitudes depend on $(x,z)$ only, are the displacement $u_i(\vect{x},t)$, the electrostatic potential $\varphi(\vect{x},t)$ and the electron-number perturbation $\Delta n(\vect{x},t)$ measured from the uniform doping $n_0$.  The polar $c$ axis is $x_3\equiv z$.  Latin indices run over $1,2,3$, a comma denotes a spatial derivative and a dot a time derivative.  The electric field is $E_i=-\partial_i\varphi$, the electron charge is $-q$ with $q>0$, and $\widetilde\mu$ is the electrochemical potential of an electron per carrier, measured in joules.

The small-strain tensor and its gradient are
\begin{equation}
\label{eq:strain}
\varepsilon_{ij}=\tfrac12\bigl(u_{i,j}+u_{j,i}\bigr),
\qquad
\eta_{ijk}\equiv\varepsilon_{ij,k}=\tfrac12\bigl(u_{i,jk}+u_{j,ik}\bigr),
\end{equation}
so that $\eta_{ijk}=\eta_{jik}$ and $\eta_{ijk}$ has the dimension of an inverse length.  The strain has six independent components and the strain gradient eighteen.

The higher-order elastic energy is taken in the single-length (isotropic-gradient) form
\begin{equation}
\label{eq:gradient_energy}
w_g=\tfrac12\,A_{ijklmn}\,\eta_{ijk}\,\eta_{lmn}
=\tfrac12\,\ell^{2}\,c_{ijlm}\,\eta_{ijk}\,\eta_{lmk},
\qquad
A_{ijklmn}=\ell^{2}\,c_{ijlm}\,\delta_{kn},
\end{equation}
in which the third and sixth slots of $A_{ijklmn}$ are the gradient directions.  This is the contraction that preserves the major symmetry $A_{ijklmn}=A_{lmnijk}$ required of a second derivative of an energy.

The units follow from the derivatives of the energy.  Stress $\sigma_{ij}=\partial w/\partial\varepsilon_{ij}$ and stiffness $c_{ijkl}$ are in pascals.  The hyperstress $\tau_{ijk}=\partial w/\partial\eta_{ijk}$ is a force per unit length (\si{\newton\per\metre}), the gradient modulus $A_{ijklmn}$ is a force (\si{\newton}), and the flexoelectric constant $f_{ijkl}=\partial D_i/\partial\eta_{jkl}$ is a charge per unit length (\si{\coulomb\per\metre}).  In the carrier sector $\Delta n$ is a number density (\si{\per\cubic\metre}), the compressibility $a=k_BT/n_0$ is an energy times a volume, and the coupling $H_{jkl}=\partial\widetilde\mu/\partial\eta_{jkl}=\partial\tau_{jkl}/\partial(\Delta n)$ is an energy times a length (\si{\joule\metre}), the one non-classical unit of the model.

\subsection{Constitutive relations and the balance laws}
\label{sec:governing_constitutive}
\label{sec:governing_balance}

The only physics added to the classical flexoelectric piezoelectric-semiconductor enthalpy is one bilinear term that couples the carrier perturbation to the strain gradient,
\begin{equation}
\label{eq:added_coupling}
\Delta\mathcal{H}=\Delta n\,H_{jkl}\,\varepsilon_{jk,l},
\end{equation}
where $H_{jkl}$ is a third-order tensor with units of energy times length.  The term belongs to a thermodynamic potential, so its two conjugates follow by differentiation and bring in no independent constants.  Differentiation with respect to the strain gradient and to the carrier density gives
\begin{equation}
\label{eq:added_hyperstress}
\Delta\tau_{jkl}=\frac{\partial\Delta\mathcal{H}}{\partial\varepsilon_{jk,l}}=H_{jkl}\,\Delta n ,
\end{equation}
\begin{equation}
\label{eq:added_mu}
\Delta\widetilde\mu=\frac{\partial\Delta\mathcal{H}}{\partial\Delta n}=H_{jkl}\,\varepsilon_{jk,l} .
\end{equation}
Only the component $H\equiv H_{zzz}$ is retained in the calculations, so \eqref{eq:added_hyperstress} adds $H\Delta n$ to $\tau_{zzz}$ and \eqref{eq:added_mu} adds $H\eta_{zzz}$ to $\widetilde\mu$.  Likewise only $f\equiv f_{3zz}$ is retained among the $6mm$-allowed flexoelectric components.  The model is therefore a reduced single-component coupling model, not the full $6mm$-invariant flexoelectric--carrier model; the remaining symmetry-allowed components are set to zero and no result reported here tests the sensitivity to that choice.

\paragraph{Relation to the thermodynamic framework.}
The closest primary work is the thermodynamically consistent theory of Qu, Pan, Zhu and Deng \cite{qu2025thermodynamic}, which derives the constitutive equations of a flexoelectric semiconductor from the free-energy imbalance and contains the drift--diffusion current, the flexoelectric and piezoelectric couplings and the higher-order stress in one framework.  Its explicit free energy, reproduced in \cite{hu2025computational}, is additive and its carrier energy contains no strain gradient, so in that form the strain gradient reaches the carriers only through the flexoelectric polarization that enters Gauss's law.  The bilinear term $\Delta n\,H_{jkl}\varepsilon_{jk,l}$ of \eqref{eq:added_coupling} is absent from that form but not from the article: its Appendix~C.1 adds a chemo-mechanical energy that is linear in the carrier concentration and in the strain gradient, and the part $M^{n}_{IJK}$ corresponds to $H_{jkl}$ here; the article neglects it in its main text and does not pose an edge-scattering problem.  The term is therefore not claimed as new, and what is new is the scattering problem it is used in.  The symbol $H$ is overloaded in the literature: in \cite{hu2025computational} it denotes the sixth-order strain-gradient elasticity tensor, written $A_{ijklmn}$ here, whereas $H_{jkl}$ is a third-order carrier--strain-gradient coefficient (J\,m) and the two are unrelated.  The incremental effect of the term on the edge partition, and the fact that it is not the origin of the backward branch, are quantified in \cref{sec:results_mechanism,sec:results_coupling}.

The electric enthalpy density of the model is
\begin{equation}
\label{eq:enthalpy}
\begin{aligned}
w={}&\tfrac12\,\varepsilon_{ij}c_{ijkl}\varepsilon_{kl}
+\tfrac12\,A_{ijklmn}\eta_{ijk}\eta_{lmn}
-\tfrac12\,E_i\kappa_{ij}E_j
-e_{kij}E_k\varepsilon_{ij}
-f_{ijkl}E_i\eta_{jkl}\\
&+\tfrac12\,a\,(\Delta n)^2
+H_{jkl}\,\Delta n\,\eta_{jkl}
-q\,\Delta n\,\varphi ,
\end{aligned}
\end{equation}
with $a=k_BT/n_0$.  The terms are, in order, the elastic, strain-gradient, dielectric, piezoelectric and flexoelectric energies of the classical flexoelectric dielectric \cite{shen2010theory}, the carrier compressibility, the added carrier--strain-gradient coupling and the electrostatic energy of the mobile charge.  The seventh term, \eqref{eq:added_coupling}, is the one added here, with the structure of the chemo-mechanical term of \cite{qu2025thermodynamic} (Appendix~C).  Differentiating \eqref{eq:enthalpy} gives the constitutive relations
\begin{align}
\sigma_{ij}&=\frac{\partial w}{\partial\varepsilon_{ij}}
=c_{ijkl}\varepsilon_{kl}-e_{kij}E_k,
\label{eq:sigma}\\
\tau_{ijk}&=\frac{\partial w}{\partial\eta_{ijk}}
=A_{ijklmn}\eta_{lmn}+H_{ijk}\Delta n-f_{lijk}E_l,
\qquad \tau_{ijk}=\tau_{jik},
\label{eq:tau}\\
D_i&=-\frac{\partial w}{\partial E_i}
=\kappa_{ij}E_j+e_{ijk}\varepsilon_{jk}+f_{ijkl}\eta_{jkl},
\label{eq:D}\\
\widetilde\mu&=\frac{\partial w}{\partial\Delta n}
=a\,\Delta n-q\varphi+H_{jkl}\eta_{jkl}.
\label{eq:mu}
\end{align}
The term-by-term derivatives and the Euler identity for $w$ are given in \cref{app:derivations}.  Every term of \eqref{eq:mu} is an energy: $k_BT\,\Delta n/n_0$ is the chemical part, $-q\varphi$ the electrostatic energy of a particle of charge $-q$, and the $H$ term the energy that the added coupling assigns to a carrier in a strain gradient.

The non-classical entries are reciprocal.  The same tensor $H_{jkl}$ appears in $\Delta n$ in the hyperstress and in $\eta_{jkl}$ in the electrochemical potential, and the same tensor $f_{ijkl}$ appears in the direct and converse flexoelectric terms of $D$ and $\tau$.  This is the Maxwell symmetry of the second derivatives of $w$, and the discretized operator inherits it: reciprocity at $\pm k_y$ is confirmed numerically to $10^{-13}$.  A single flexoelectric component, $f\equiv f_{3zz}$, is retained, so that $D_z=e_{zjk}\varepsilon_{jk}+\kappa_{zz}E_z+f\,\eta_{zzz}$ and $\tau_{zzz}=A_{zzz\,lmn}\eta_{lmn}-fE_z$.  The surface flexoelectric effect \cite{shen2010theory} and electric-field-gradient flexoelectric terms are not included.

For the $6mm$ class with the polar axis $x_3\equiv z$, the tensors reduce to five elastic constants, three piezoelectric constants and two permittivities.  In the Voigt ordering $(\varepsilon_{xx},\varepsilon_{yy},\varepsilon_{zz},\gamma_{yz},\gamma_{xz},\gamma_{xy})$, with $\gamma_{ij}=2\varepsilon_{ij}$, the elastic and piezoelectric matrices of ZnO are
\begin{equation}
\label{eq:CZnO}
\mathsf C=
\begin{pmatrix}
c_{11}&c_{12}&c_{13}&0&0&0\\
c_{12}&c_{11}&c_{13}&0&0&0\\
c_{13}&c_{13}&c_{33}&0&0&0\\
0&0&0&c_{44}&0&0\\
0&0&0&0&c_{44}&0\\
0&0&0&0&0&c_{66}
\end{pmatrix},
\qquad
\begin{aligned}
c_{11}&=\SI{210}{\giga\pascal},& c_{33}&=\SI{221}{\giga\pascal},\\
c_{13}&=\SI{105}{\giga\pascal},& c_{12}&=\SI{121}{\giga\pascal},\\
c_{44}&=\SI{43}{\giga\pascal},& c_{66}&=\tfrac12(c_{11}-c_{12}),
\end{aligned}
\end{equation}
{\small\begin{equation}
\label{eq:ZnOe}
\mathsf e=
\begin{pmatrix}
0&0&0&0&e_{15}&0\\
0&0&0&e_{15}&0&0\\
e_{31}&e_{31}&e_{33}&0&0&0
\end{pmatrix},
\qquad
e_{33}=\SI{1.32}{\coulomb\per\square\metre},\;
e_{31}=\SI{-0.57}{\coulomb\per\square\metre},\;
e_{15}=\SI{-0.48}{\coulomb\per\square\metre},
\end{equation}}
and the permittivity is diagonal, $\kappa_{11}=\kappa_{22}$, $\kappa_{33}=\SI{8.85e-11}{\farad\per\metre}$.  The component relations used are
\begin{equation}
\label{eq:sigmaZnO}
\begin{aligned}
\sigma_{xx}&=c_{11}\varepsilon_{xx}+c_{12}\varepsilon_{yy}
+c_{13}\varepsilon_{zz}-e_{31}E_z,\\
\sigma_{yy}&=c_{12}\varepsilon_{xx}+c_{11}\varepsilon_{yy}
+c_{13}\varepsilon_{zz}-e_{31}E_z,\\
\sigma_{zz}&=c_{13}(\varepsilon_{xx}+\varepsilon_{yy})
+c_{33}\varepsilon_{zz}-e_{33}E_z,\\
\sigma_{yz}&=2c_{44}\varepsilon_{yz}-e_{15}E_y,\qquad
\sigma_{xz}=2c_{44}\varepsilon_{xz}-e_{15}E_x,\qquad
\sigma_{xy}=2c_{66}\varepsilon_{xy},
\end{aligned}
\end{equation}
\begin{equation}
\label{eq:tauZnO}
\tau_{zzz}=\ell^{2}\Bigl[c_{13}\bigl(\eta_{xxz}+\eta_{yyz}\bigr)
+c_{33}\,\eta_{zzz}\Bigr]-f\,E_z+H\,\Delta n ,
\end{equation}
\begin{equation}
\label{eq:DZnO}
\begin{aligned}
D_x&=\kappa_{11}E_x+2e_{15}\varepsilon_{xz},\qquad
D_y=\kappa_{11}E_y+2e_{15}\varepsilon_{yz},\\
D_z&=\kappa_{33}E_z+e_{31}(\varepsilon_{xx}+\varepsilon_{yy})
+e_{33}\varepsilon_{zz}+f\,\eta_{zzz}.
\end{aligned}
\end{equation}
The field equations are of fourth order in the mechanical sector and of second order in the electrical and carrier sectors, which is why the plate supports four wave families per branch.

The field equations are the Euler--Lagrange equations of the enthalpy.  The variations $\delta\varepsilon_{ij}$ and $\delta\eta_{ijk}=\tfrac12(\delta u_{i,jk}+\delta u_{j,ik})$ follow from \eqref{eq:strain}, and because $\tau_{ijk}=\tau_{jik}$ one has $\tau_{ijk}\,\delta\eta_{ijk}=\tau_{ijk}\,\delta u_{i,jk}$.  One integration by parts on the stress and two on the hyperstress move all derivatives onto the fields, and the surface terms are the traction and the double traction,
\begin{equation}
\label{eq:ibp_hyper}
\int_V\tau_{ijk}\,\delta\eta_{ijk}\,\dd V
=\int_V\partial_k\partial_j\tau_{ijk}\,\delta u_i\,\dd V
-\oint_{\partial V}n_j\partial_k\tau_{ijk}\,\delta u_i\,\dd S
+\oint_{\partial V}n_k\tau_{ijk}\,\delta u_{i,j}\,\dd S .
\end{equation}
With the inertial term, the mechanical balance is
\begin{equation}
\label{eq:momentum}
\rho\,\ddot u_j=\partial_i\sigma_{ij}-\partial_i\partial_k\tau_{jik},
\end{equation}
where the double divergence of the hyperstress enters with a minus sign because the gradient energy acts against the elastic restoring force.  With the effective stress
\begin{equation}
\label{eq:effective}
\Sigma_{ij}=\sigma_{ij}-\partial_k\tau_{ijk},
\end{equation}
the balance reads $\rho\ddot u_j=\partial_i\Sigma_{ij}$ and the traction on a surface is $t_j=n_i\Sigma_{ij}$.  The expanded out-of-plane component of \eqref{eq:momentum} and of Gauss's law are given in \cref{app:derivations}.

Variation with respect to $\varphi$ gives Gauss's law,
\begin{equation}
\label{eq:gauss}
\partial_iD_i=-q\,\Delta n .
\end{equation}
The carrier equation is not a variational equation of the enthalpy: the derivative with respect to $\Delta n$ only defines $\widetilde\mu$ in \eqref{eq:mu}, and the continuity statement follows from conservation of carrier number with the drift--diffusion flux,
\begin{equation}
\label{eq:continuity}
\partial_t(\Delta n)+\partial_iJ_i=0,
\qquad
J_i=-\frac{D\,n_0}{k_BT}\,\partial_i\widetilde\mu ,
\end{equation}
where $J_i$ is the electron number flux and $D=\mu_ek_BT/q$ the diffusivity (Einstein relation).  With $\partial_t\to-\ii\omega$ this gives $\Delta n=\ii D n_0\nabla^{2}\widetilde\mu/(\omega k_BT)$, and the factor $\ii$ reappears as $+\ii\alpha_{\mathrm d}/\omega$ in the weak form \eqref{eq:weak_carrier}; reversing its sign would turn the carrier sector from a sink of energy into a source.  The charge-current density is
\begin{equation}
\label{eq:current}
j_i=-qJ_i=\frac{qDn_0}{k_BT}\,\partial_i\widetilde\mu
=qD\,\partial_i\Delta n-q\mu_e n_0\,\partial_i\varphi
+\frac{qDn_0}{k_BT}H_{jkl}\,\partial_i\eta_{jkl},
\end{equation}
whose three terms are the diffusion current, the drift current and the carrier--strain-gradient contribution.  The last term is the transport signature of \eqref{eq:added_coupling}: a strain-gradient field drives a current even at zero applied field.

\subsection{The electric model, the boundary conditions and the energy flux}
\label{sec:governing_electric}
\label{sec:governing_bc}
\label{sec:governing_energy}

Equations \eqref{eq:gauss}--\eqref{eq:continuity} are electrostatic: no magnetic induction, displacement current or electromagnetic radiation mode is retained.  This is justified by the ratios that control it.  At $f=\SI{16.68}{\giga\hertz}$ and with $\kappa_{33}/\varepsilon_0=10$ the electromagnetic wavenumber is $k_{\mathrm{EM}}=\SI{1.106e3}{\per\metre}$, so $k_{\mathrm{EM}}d=2.2\times10^{-4}$ for $d=\SI{200}{\nano\metre}$ and the retardation correction is of order $(k_{\mathrm{EM}}d)^2=4.9\times10^{-8}$.  The phase speed gives the same result: $v_p/c=9.4\times10^{-5}$, so the magnetic energy is smaller than the electric energy by $(v_p/c)^2\approx8.8\times10^{-9}$.

The carriers are different.  With $\sigma_e=q\mu_en_0=\SI{5.6e-3}{\siemens\per\metre}$ the dielectric relaxation time is $\tau_d=\kappa_{33}/\sigma_e=\SI{1.58e-8}{\second}$ against a wave period of $\SI{6.0e-11}{\second}$, so $\omega\tau_d=1.65\times10^{3}$ and the carriers cannot follow the field quasistatically.  The model therefore keeps the full continuity equation \eqref{eq:continuity}; what is neglected is electromagnetic retardation, not charge dynamics.  The Debye length is $L_D=\sqrt{\kappa_{33}k_BT/(q^2n_0)}=\SI{2.68}{\micro\metre}$, so $d/L_D=0.075$ and the carrier perturbation is not confined to a surface layer.

The boundary conditions follow from the surface terms of the variational principle.

\paragraph{Plate faces $z=\pm d/2$.}  The faces are free surfaces, so the generalized traction of \eqref{eq:generalised_traction} (with face normal $z$) and the double traction vanish; the electrons cannot leave the crystal; and the electrical condition is grounded or open circuit,
\begin{equation}
\label{eq:bc_faces_mech}
\bar\sigma_{iz}=0,\qquad \tau_{izz}=0,\qquad z=\pm d/2,
\end{equation}
\begin{equation}
\label{eq:bc_faces_carrier}
J_z=0,\qquad z=\pm d/2,
\end{equation}
\begin{equation}
\label{eq:bc_grounded}
\varphi=0\quad\text{(grounded)},\qquad
D_z=0\quad\text{(open circuit)},\qquad z=\pm d/2 .
\end{equation}
The grounded condition models perfectly conducting electrodes shorted to each other; it is essential and is imposed on the basis functions.  The open-circuit condition models a bare surface with no free charge; it is natural and is imposed weakly, with no potential mode removed.  A potential that is constant across the thickness is a gauge only in the sector $k_x=k_y=0$, which is not used here; for $(k_x,k_y)\neq(0,0)$ it gives physical in-plane fields.  The open-circuit results are checked against the independent $C^{1}$ discretization, which tests the discretization and not the physics.

\paragraph{Reflecting edge $x=0$.}  The edge is a free surface in vacuum and no datum is prescribed on it, so all four conditions are natural, each acting on the test functions of the field it is conjugate to,
\begin{equation}
\label{eq:bc_edge}
\bar\sigma_{ix}=0,\qquad \tau_{ixx}=0,\qquad D_x=0,\qquad J_x=0,\qquad x=0 ,
\end{equation}
where
\begin{equation}
\label{eq:generalised_traction}
\bar\sigma_{i}(\vect n)=n_k\underbrace{\bigl(\sigma_{ik}-\partial_l\tau_{ikl}\bigr)}_{\Sigma_{ik}}
-\partial_j^{s}\bigl(n_k\tau_{ijk}\bigr),
\end{equation}
is the generalized traction on a surface of unit normal $\vect n$, with $\partial^s_j$ the derivative tangential to that surface: the effective stress minus the surface divergence of the hyperstress.  For the edge $\vect n=\vect e_x$ it is written $\bar\sigma_{ix}$, and for the faces $\vect n=\pm\vect e_z$ it gives $\pm\bar\sigma_{iz}$, with $j=x,y$ tangential there.  The free index is the first index of the hyperstress and its gradient index is contracted with the normal, so the edge term is $\partial_j^s\tau_{ijx}$ and not $\partial_j^s\tau_{ixj}$; a wrong index changes the result.  With $n=e_x$ and a tangential wavenumber $k_y$, $\bar\sigma_{jx}=\Sigma_{jx}-\ii k_y\tau_{jyx}-\partial_z\tau_{jzx}$.  The second entry of \eqref{eq:bc_edge} is the double traction, conjugate to the normal derivative of the displacement on the edge.  In this model $\tau_{ijx}=\ell^{2}c_{ijkl}\eta_{klx}$, so the double traction is present if and only if $\ell\neq0$; no result reported here takes the limit $\ell=0$ with $f\neq0$.  The potential is free on the edge, so $D_x=0$ is its natural condition.  The surviving terms are set out case by case in \cref{app:derivations}.

\paragraph{Far field $x\to-\infty$.}  The solution is the prescribed incident guided mode plus outgoing branches; no other wave arrives from $x=-\infty$.

\paragraph{The corner.}  Where the edge meets a face, the surface divergence in \eqref{eq:generalised_traction} has a boundary, and the integration by parts that produced the traction also produces a line force along the corner.  With $\nu$ the co-normal of the edge boundary inside the face, its density per unit length is
\begin{equation}
\label{eq:corner_force}
f_i=\nu_j\,n_k\,\tau_{jik},
\end{equation}
summed over the two free faces that meet along the corner.  For this geometry the contributions are the mixed components $\tau_{i31}$ and $\tau_{i13}$, so $f_i=\tau_{i31}+\tau_{i13}$.  The relative sign and the index placement are fixed, together with \eqref{eq:generalised_traction} and the double traction $m_i=\tau_{i11}$, by the identity
\begin{equation}
\label{eq:corner_identity}
\bigl\langle T_{\!u}-t\bigr\rangle
=f_+\,b_0(+\tfrac12)-f_-\,b_0(-\tfrac12),
\qquad
\bigl\langle R-m\bigr\rangle=0 ,
\end{equation}
between the derived fields and the discrete flux rows $T_{\!u}$, $R$ of the canonical state, with $\langle\cdot\rangle$ the projection on the displacement test basis $b_0$.  The identity holds to $1.4\times10^{-13}$ relative for the incident guided mode at all degrees from 10 to 28, and it is violated at the level of $1.7\times10^{1}$ if the corner terms are dropped, at $5.7\times10^{-1}$ if only $\tau_{i31}$ is kept, and at $1.1\times10^{0}$ if the two are subtracted instead of added, so the corner line force is an essential part of the natural data.  It is the one natural datum that a polynomial Galerkin basis cannot annihilate pointwise; its magnitude and decay are quantified in the conclusions.

The power flowing in the $x$ direction, per unit width and integrated over the thickness, is
\begin{equation}
\label{eq:Px_int}
\mathcal{P}_x=\tfrac12\Real\int_{-d/2}^{d/2}
\Bigl[-\Sigma_{xi}\dot u_i^{*}-\tau_{jkx}\dot\varepsilon_{jk}^{*}
+\varphi\dot D_x^{*}+\widetilde\mu\,J_x^{*}\Bigr]\dd z
-\tfrac12\Real\Bigl[\tau_{ixz}\dot u_i^{*}\Bigr]_{z=-d/2}^{z=d/2} ,
\end{equation}
where $\Sigma_{xi}$ is the effective stress \eqref{eq:effective} and $J_x$ is the $x$ component of the electron-number flux of \eqref{eq:continuity}.  The four integrands are the elastic, hyperelastic, dielectric and chemical contributions, and the last term is the line work of the double stress on the two faces.  It does not vanish for a guided mode and it is needed for the flux to balance the loss: without it the single-mode balance $\lvert2\,\Imag k_x\,\mathcal P_x\rvert=P_{\text{diss}}$ is violated at the tens-of-percent level, and with it the balance closes to $3\times10^{-9}$ or better.  The flux $\mathcal P_x$, and not the magnitude of the wavevector, decides whether a mode is backward: a mode is backward when $\mathcal{P}_x<0$ while it is outgoing at the edge.

For a scattering solution the identity
\begin{equation}
\label{eq:energy_identity}
\underbrace{P_{\text{inc}}}_{\text{incident}}
+\underbrace{\sum_{j}P_j}_{\text{self-powers}}
+\underbrace{P_{\text{cross}}}_{\text{interference}}
=\underbrace{P_{\text{edge}}+P_{\text{diss}}}_{\text{edge flux plus loss}}
\end{equation}
holds exactly, with $P_j$ the signed self-power of outgoing mode $j$, $P_{\text{cross}}$ the bilinear interference between modes that are not orthogonal in the power metric, and
\begin{equation*}
P_{\text{diss}}=\tfrac12\alpha_{\mathrm{d}}\int_{-L}^{0}\!\!\int
\Bigl[|\partial_x\widetilde\mu|^{2}+k_y^2|\widetilde\mu|^2+|\partial_z\widetilde\mu|^2\Bigr]\dd z\,\dd x ,
\qquad \alpha_{\mathrm{d}}=\frac{D}{d\,c_s},
\end{equation*}
written in scaled variables, with $\widetilde\mu$ the total electrochemical-potential field divided by $k_BT$.  The integrand is nonnegative, so $P_{\text{diss}}\ge0$, and $\alpha_{\mathrm d}$ is the ratio of the diffusion rate to the elastic rate, a property of the material and thickness and not of a mode (the spatial attenuation of a mode is $\alpha=2|\Imag k_x|$).  Only the propagating members of $\{P_j\}$ transport energy to infinity; the evanescent members are reactive at the plane and $P_{\text{cross}}$ is an interference correction.  The residual of \eqref{eq:energy_identity} is evaluated for each reported solution; its values at the operating point for the five coupling levels are plotted in \cref{fig:fig05}(c).

\subsection{Plane-wave reduction and the characteristic polynomial}
\label{sec:governing_reduction}

The plate is invariant along $y$ and in time, so a single tangential wavenumber and
frequency are conserved and every field takes the form
$f=\widehat f(z)\,\ee^{\ii(k_xx+k_yy-\omega t)}$, for which
\begin{equation}
\label{eq:substitution}
\partial_t\to-\ii\omega,\qquad
\partial_x\to\ii k_x,\qquad
\partial_y\to\ii k_y,\qquad
\partial_z\to\frac{\dd}{\dd z}\equiv{}' .
\end{equation}
The strain, strain-gradient and field components follow by applying
\eqref{eq:substitution} to \eqref{eq:strain}; for example
$\widehat\varepsilon_{xx}=\ii k_x\widehat u_1$,
$\widehat\eta_{xxx}=-k_x^{2}\widehat u_1$, $\widehat\eta_{zzz}=\widehat u_3''$ and
$\widehat E_x=-\ii k_x\widehat\varphi$.  The continuity equation
\eqref{eq:continuity} becomes algebraic, so the carrier amplitude is carried by the
same amplitude vector as the displacements and has no independent time evolution.

The dependence on $k_x$ is fixed by derivative counting.  The stress $\sigma$ is linear in
$\varepsilon$ and $E$, each containing one power of $k_x$, and enters the
mechanical balance through one further derivative, so that term is at most quadratic
in $k_x$.  The hyperstress $\tau$ is linear in $\eta$, which contains two powers, and
enters through two derivatives, so that term is at most quartic.  Counted on the assembled
operator, each side of the weak form carries at most two powers (one from $\varepsilon$ or $E$, two
from $\eta$, none from $\Delta n$ or $\widetilde\mu$), so no power above four arises, and four arises
only from the $\eta$--$\eta$ gradient energy.  The retained couplings add nothing at that order: $H\,\Delta n\,\eta_{zzz}$ pairs a
zero-power factor with $\eta_{zzz}=\widehat u_3''$, which contains no $k_x$, and $f_{3zz}E_z\eta_{zzz}$ involves only $z$ derivatives.  For a
vector $\vect q$ of thickness degrees of freedom the reduced problem is therefore a
polynomial of degree four in $k_x$,
\begin{equation}
\label{eq:quartic}
\Bigl[\mathsf A_{00}
+\ii k_x\bigl(\mathsf A_{01}-\mathsf A_{10}\bigr)
+k_x^{2}\mathsf A_{11}
+\ii k_x^{3}\bigl(\mathsf A_{12}-\mathsf A_{21}\bigr)
+k_x^{4}\,\mathsf A_{22}\Bigr]\vect q=\vect 0 ,
\end{equation}
where $\mathsf A_{rs}$ collects the terms containing $\partial_x^{r}$ on the left test
function and $\partial_x^{s}$ on the right one (their construction is given with the
discretization), and the $k_x^{2}$ coefficient reduces to $\mathsf A_{11}$ because
$\mathsf A_{02}+\mathsf A_{20}$ does not arise from this enthalpy.  The odd coefficients
are Hermitian, being $\ii$ times antisymmetric differences; the even ones carry the
non-Hermitian carrier transport form.  The degree is exactly four only for $\ell\neq0$:
for the retained $f_{3zz}$ and $H_{zzz}$, $\mathsf A_{22}$ and $\mathsf A_{12}$ come from the strain-gradient energy alone, so
for $\ell=0$ (with $f$ and $H$ removed) $\mathsf A_{22}=\mathsf A_{12}=0$ and
\eqref{eq:quartic} collapses to the quadratic problem of the classical
piezoelectric-semiconductor plate.  The fourth-order structure is thus a property
of the strain-gradient energy, not of the carrier physics.

Read by sectors and not by powers of $k_x$, the same operator has a $3\times3$ block
structure.  With $\mathsf B^{,l}$ the strain operator differentiated once more along $l=x,y,z$ and unknowns ordered as
$(\widehat{\vect u},\widehat{\boldsymbol\varphi},\widetilde{\boldsymbol\mu})$,
\begin{equation}
\label{eq:sector_blocks}
\mathsf A(k_x,\omega)=
\begin{pmatrix}
\mathsf A_{uu} & \mathsf A_{u\varphi} & \mathsf A_{u\mu}\\[2pt]
\mathsf A_{\varphi u} & \mathsf A_{\varphi\varphi} & \mathsf A_{\varphi\mu}\\[2pt]
\mathsf A_{\mu u} & \mathsf A_{\mu\varphi} & \mathsf A_{\mu\mu}
\end{pmatrix},
\qquad
\begin{aligned}
\mathsf A_{uu}&=\langle\mathsf B,\mathsf C\mathsf B\rangle
+\lambda^{2}\sum_{l}\langle\mathsf B^{,l},\mathsf C\mathsf B^{,l}\rangle
-\omega^{2}\langle\mathsf U,\mathsf U\rangle,\\
\mathsf A_{u\varphi}&=\langle\mathsf B,\mathsf e^{\mathsf T}\mathsf G\rangle
+f\langle\partial_\zeta\mathsf B^{(2)},\mathsf G^{(2)}\rangle,\\
\mathsf A_{\varphi\varphi}&=-\langle\mathsf G,\boldsymbol\kappa\mathsf G\rangle,\qquad
\mathsf A_{\varphi\mu}=-Q\langle\mathsf G^{(0)},\mathsf B\rangle,\\
\mathsf A_{\mu\mu}&=-\langle\mathsf W,\mathsf W\rangle
-\frac{\ii\alpha_{\mathrm d}}{\omega}\langle\mathsf M\mathsf G,\mathsf M\mathsf G\rangle,
\end{aligned}
\end{equation}
where the carrier form $-\langle\mathsf W,\mathsf W\rangle$ of \eqref{eq:block}, with
$\mathsf W=(-h\,b_j''\text{ in the }\widehat u_3\text{ slots},\;Q\psi_j,\;b_j)$ and the
prime the derivative with respect to $z/d$, also contributes
$-h^{2}\langle b'',b''\rangle$ to $\mathsf A_{uu}$, $+hQ\langle b'',\psi\rangle$ to
$\mathsf A_{u\varphi}$, $+h\langle b'',b\rangle$ to $\mathsf A_{u\mu}$ and
$-Q^{2}\langle\psi,\psi\rangle$ to $\mathsf A_{\varphi\varphi}$.  The remaining
blocks are the Hermitian transposes of those written, being the same derivatives of the
same enthalpy taken in the other order.  The symmetric part of \eqref{eq:sector_blocks}
is the Hessian of the discretized enthalpy; the carrier transport form is the only
antisymmetric contribution and the entire source of loss.

\subsection{Non-dimensionalization, limits and admissibility}
\label{sec:governing_nondim}
\label{sec:governing_limits}
\label{sec:governing_stability}

The calculation uses lengths scaled by $d$, stresses by $c_{44}$, time by
$d/c_s$ with $c_s=\sqrt{c_{44}/\rho}$, carrier densities by $n_0$ and energies
per carrier by $k_BT$.  The carrier amplitude of \eqref{eq:continuity} is
measured so that the compressibility term of \eqref{eq:enthalpy} is
$\tfrac12\int\widehat N^{2}$ in the scaled units used by \eqref{eq:block}, which
is what fixes the coefficient of $\widehat N^{2}$ there.  Two derived scales appear repeatedly,
\begin{equation*}
a=\frac{k_BT}{n_0},\qquad L_D=\sqrt{\frac{\kappa_{33}k_BT}{q^2n_0}},
\end{equation*}
and the model is governed by the five dimensionless numbers of
\cref{tab:groups}.  Each is an exact ratio of two physical lengths, times or
energies, and each follows directly from the parameters listed below.

\begin{table}[t]
\centering
\caption{Dimensionless groups and their values at the operating point; $T=\SI{301.7}{\kelvin}$ follows from the Einstein relation; $h$ is given at the largest coupling examined, $|H|=\SI{2.0e-23}{\joule\metre}$.}
\label{tab:groups}
\begin{tabular}{@{}llll@{}}
\toprule
Group & Definition & Value & Meaning \\
\midrule
$Q$ & $d/L_D$ & $7.46\times10^{-2}$ & thickness over Debye length \\
$\alpha_{\mathrm{d}}$ & $D/(d\,c_s)$ & $8.28\times10^{-1}$ & diffusion rate over elastic rate \\
$\lambda$ & $\ell/d$ & $2.50\times10^{-1}$ & gradient length over thickness \\
$f$ & $f_{3zz}/(d\sqrt{c_{44}\kappa_{33}})$ & $2.56\times10^{-3}$ & flexoelectric coupling \\
$h$ & $H/(d\sqrt{a\,c_{44}})$ & $\pm1.06\times10^{-2}$ & carrier--gradient coupling (at $|H|_{\max}$) \\
\bottomrule
\end{tabular}
\end{table}

\paragraph{Parameters and their provenance.}\label{sec:params}
The classical constants are literature values ($c_{11}=210$, $c_{33}=221$, $c_{44}=43\,\si{\giga\pascal}$; $e_{33}=1.32$, $e_{31}=-0.57$, $e_{15}=-0.48\,\si{\coulomb\per\square\metre}$; $\kappa_{33}=8.85\times10^{-11}\,\si{\farad\per\metre}$, from Table~1 of \cite{zhang2025pscplate}).  The transport parameters are representative values for n-type ZnO: $\mu_e=0.0175\,\si{\square\metre\per\volt\per\second}$ and $D=4.55\times10^{-4}\,\si{\square\metre\per\second}$, equal to $\mu_ek_BT/q$ at $T=\SI{301.7}{\kelvin}$ and within the $120$--$\SI{205}{\square\centi\metre\per\volt\per\second}$ compiled for nominally undoped bulk and film ZnO in Table~VII of \cite{ozgur2005zno}; the density $n_0=1$--$3\times10^{18}\,\si{\per\cubic\metre}$ is a modeling scenario below the $8\times10^{19}$--$5\times10^{23}\,\si{\per\cubic\metre}$ of that table, and $d=\SI{200}{\nano\metre}$ is a reference thickness (sweep $120$--$\SI{320}{\nano\metre}$).  The strain-gradient length $\ell=\SI{50}{\nano\metre}$, the flexoelectric coefficient $f_{3zz}=\SI{1}{\nano\coulomb\per\metre}$ (in the nC/m range reported for ZnO \cite{hu2025computational,zhao2021effects}; it contributes \SI{-0.01}{\percent} to $R_B$) and the coupling $H_{333}$ are \emph{assumed}; the same form $g=\ell^2c$ is used in \cite{qu2023drag} with \SI{20}{\nano\metre} for silicon.  For $H_{jkl}$ no published value is available: the numbers used, $|H|\le\SI{2.0e-23}{\joule\metre}=0.022\,H_{\text{bound}}$, are fixed by the positivity bound \eqref{eq:schur}, which leaves a margin of a factor of forty-five in $H$.

The groups follow from the listed parameters by algebra, so their verification is a comparison of numbers and not a fit.  The elastic speed is
$c_s=\sqrt{c_{44}/\rho}=\SI{2.7466}{\kilo\metre\per\second}$, giving the frequency unit
$c_s/(2\pi d)=\SI{2.1857}{\giga\hertz}$ at $d=\SI{200}{\nano\metre}$; the operating point
$\omega=7.6334\,c_s/d$ is then $f=\SI{16.684}{\giga\hertz}$, the value quoted in the abstract.
With $a=k_BT/n_0=\SI{2.0828e-39}{\joule\metre^{3}}$ (the temperature
$T=\SI{301.72}{\kelvin}$ is implied by the Einstein relation for the diffusivity--mobility
pair, not an independent input), the Debye length is
$L_D=\SI{2.6797}{\micro\metre}$ and $Q=d/L_D=0.074635$.  The plate is thus thinner than
$L_D/d=13.40$ times a screening length, the carrier perturbation is not confined to a
surface layer, and the screening enters \eqref{eq:block} only through the small coefficient
$Q^{2}=5.57\times10^{-3}$.  The transport group is
$\alpha_{\mathrm d}=D/(d\,c_s)=0.8283$, of order unity: the diffusion time across the
thickness, $d^{2}/D=\SI{8.8e-11}{\second}$, is comparable to the elastic transit time
$2d/c_s=\SI{1.5e-10}{\second}$, so the carrier and elastic sectors are not separated by a
small parameter and the carrier cannot be condensed away.  Finally $\lambda=\ell/d=0.25$
and $h=H/H_{\mathrm{ref}}$ with
$H_{\mathrm{ref}}=d\sqrt{a\,c_{44}}=\SI{1.8927e-21}{\joule\metre}$, which has the units of $H$
required by \eqref{eq:enthalpy}; the tabulated extreme $|H|=\SI{2.0e-23}{\joule\metre}$ is
$h=\pm1.06\times10^{-2}$, a one-percent perturbation of the carrier sector.

The model of \eqref{eq:enthalpy} contains four sectors that can be switched off
independently, and the limits they generate are the ones against which the
calculation is read.

\emph{Classical piezoelectric plate.}  Setting $\ell=f=H=0$ and removing the
carrier sector ($\Delta n\equiv0$) leaves the classical piezoelectric
formulation, with a polynomial in $k_x$ of degree two instead of four.

\emph{Strain-gradient dielectric.}  With $H=0$ and no carriers the enthalpy is
the flexoelectric dielectric enthalpy of \cite{shen2010theory} completed by the
gradient energy \eqref{eq:gradient_energy}.  In this limit the fundamental
shear-horizontal branch has the closed form
$\Omega^{2}=(c_{66}/c_{44})K^{2}(1+\lambda^{2}K^{2})$ in the scaled units of this subsection, which is used as one of the checks.

\emph{Piezoelectric semiconductor without the gradient sector.}  Setting
$\ell=f=H=0$ and retaining the carriers leaves a drift--diffusion plate whose
sagittal spectrum is examined in \cref{sec:results_mechanism}: at the frequencies of the controls it
contains no outgoing mode with a positive axial wavenumber and a negative axial
power apart from the specular reflection.  The plate nevertheless has a backward
interval of its own; the gradient length places the fold of that interval at the
reported frequency.

\emph{Uncoupled carriers.}  With $H=0$ and the strain-gradient sector retained,
the carrier dynamics still enter through Gauss's law and through the continuity
equation \eqref{eq:continuity}, so the plate remains a piezoelectric
semiconductor; only the direct carrier--gradient enthalpy of
\eqref{eq:added_coupling} is removed.  This is the limit in which the controls are computed, and it is the limit that shows the
partition to be a property of the branch and not of the added term.

The enthalpy \eqref{eq:enthalpy} is a quadratic form in
$(\varepsilon,\eta,E,\Delta n)$.  It is a saddle in the electric field, because the
term $-\tfrac12E_i\kappa_{ij}E_j$ is negative, so the full form is not positive definite and
none of the statements below concerns it.  The block in $(\varepsilon,\eta,\Delta n)$ at
fixed $E$ is positive definite only if the carrier and strain-gradient parts are not over-coupled.  Eliminating $\Delta n$ blockwise
gives the Schur condition
\begin{equation}
\label{eq:schur}
s\equiv\frac{H^2\,(c^{-1})_{33}}{\ell^2\,a}<1 ,
\end{equation}
where $(c^{-1})_{33}$ is the $33$ entry of the inverse of the full elastic
matrix, i.e.\ the reduced modulus
$c^{*}=c_{33}-2c_{13}^2/(c_{11}+c_{12})=\SI{154.4}{\giga\pascal}$ and not $c_{33}$
itself.  The bound is attained at
$H_{\text{bound}}=\ell\sqrt{a\,c^{*}}=\SI{8.97e-22}{\joule\metre}$, which is the
value at which the coupled stiffness of the plate ceases to be positive for
every strain-gradient field.  Every calculation reported below is confined to \eqref{eq:schur}.  All reported solutions have
$|H|\le\SI{2.0e-23}{\joule\metre}=0.022\,H_{\text{bound}}$, i.e.\ $s\le
5\times10^{-4}$: the $(\varepsilon,\eta,\Delta n)$ block is positive definite with a margin of more than two
thousand, and no reported result lies near a constitutive instability.

The figures normalize the coupling by
$H^{*}=H_{\text{bound}}/\sqrt5=\SI{4.01e-22}{\joule\metre}$, which is the value
at which $s=0.2$; the normalization is a plotting convention with no physical
content beyond fixing the swept interval $|H/H^{*}|\le0.05$, that is
$|H|\le2.0\times10^{-23}\,\si{\joule\metre}$, the tabulated extreme.  The
separate symbol $H_{\mathrm{ref}}=d\sqrt{a\,c_{44}}=\SI{1.89e-21}{%
\joule\metre}$ used in \cref{tab:groups} is the thickness-based
reference of \eqref{eq:enthalpy}; it is not the plotting normalization, and the two differ.  No strain-gradient elasticity measurement for
ZnO fixes any of $\ell$, $f$ or $H$, so the reported intervals are modeling
choices and have no experimental uncertainty attached; the numerical uncertainty of a
reported $R_B$ is its refinement spread, below $3\times10^{-5}$ in relative terms.
